Gehören umgekehrt $R_\ma$ und $R_{\bar\beta}$ derselben Darstellungsklasse an, wird also $R_{\bar\beta}=P^{-1}R_\ma P$ --- wobei die $a_i$ eine Basis von $\ma$, die $\bar\beta_i$ eine Basis von $\mb$ bilden ---, so wird auch durch
\[
(\beta_1,\ldots,\beta_s)=(\bar\beta_1,\ldots,\bar\beta_s)P^{-1}
\]
eine Basis von $\mb$ gebildet, und es wird $R_\beta=PR_{\bar\beta}P^{-1}=R_\ma$. Ordnet man also jedem Element $a=c_1a_1+\cdots+c_sa_s$ aus $\ma$ das Element $\beta=c_1\beta_1+\cdots+c_s\beta_s$ aus $\mb$ zu, so ist die Zuordnung eineindeutig, und Summe und Differenz entsprechen sich. Aus $R_\ma=R_\beta$ folgt aber weiter, daß $\gamma a_i$ und $\gamma\beta_i$ sich entsprechen, da sie sich durch dieselben Linearformen in den $a_i$ bzw. $\beta_i$ ausdrücken lassen; somit entsprechen sich allgemein $\gamma a$ und $\gamma\beta$; die Ideale $\ma$ und $\mb$ gehören derselben Idealklasse an. Als \emph{Hauptklasse} wird die Klasse des Einheitsideals bezeichnet; jede Basis von $\mR$ ergibt also als Darstellungsklasse die Hauptklasse.